\section{Deep Learning-Based Radiomics Model for PD-L1 Prediction}

Li et al.\ (2024) developed a deep learning radiomics (DLR) model to predict PD-L1 expression in non-small cell lung cancer (NSCLC) using multimodal PET/CT imaging data~\cite{Li2024-np}. The study included 136 patients diagnosed with NSCLC between January 2021 and September 2022, who underwent pre-treatment PET/CT scans. The dataset was divided into a training cohort (n = 95) and a validation cohort (n = 41), following a 7:3 ratio. A total of 2,446 handcrafted radiomic features and 4,096 deep learning features were extracted from both PET and CT modalities. Statistical and machine learning methods were used to select and combine features relevant to PD-L1 status. Specifically, the Mann-Whitney U test and Spearman correlation analysis were applied for feature screening, followed by least absolute shrinkage and selection operator (LASSO) regression to identify the most predictive features. Three models were constructed: a radiomics-only model, a deep learning-only model, and a fusion model that combined both feature sets.
Model performance was assessed using receiver operating characteristic (ROC) analysis. In the training cohort, the fusion model achieved the highest performance, with an AUC = 0.954 (95\% CI: 0.890-0.986), outperforming both the radiomics model (AUC = 0.829) and the deep learning model (AUC = 0.935). Similar trends were observed in the validation cohort, where the fusion model reached an AUC = 0.910 (95\% CI: 0.779-0.977), again exceeding the performance of the radiomics (AUC = 0.785) and deep learning (AUC = 0.867) models. Sensitivity, specificity, and accuracy were also reported, confirming the robustness of the fusion approach. 

Although the methodology strongly implies the use of deep learning and radiomics analysis frameworks, the specific software libraries or programming environments employed (e.g., PyTorch, TensorFlow, PyRadiomics) were not disclosed in the publication. Furthermore, no public GitHub repository or source code is provided, and the authors state that the datasets are available upon reasonable request. Overall, this study demonstrates that integrating handcrafted radiomic features with deep learning representations from PET/CT scans can significantly improve the non-invasive prediction of PD-L1 expression in NSCLC patients.